\subsection{回転面の面積と線重心}

回転軸と同じ平面内にある曲線を考える。この曲線を回したときにできる面の面積では，面積重心ではなく線重心が現れる。その理由を，回転面の小さな部分から考える。軸の片側にあり，軸からの距離が $r_1,r_2$ の二点を結ぶ短い線分を回すと，一般には円すい台の側面になる。軸に平行な線分は円筒面，垂直な線分は円環または円板を作る。いずれの場合も，線分の長さを $\Delta s$ とすると面積は
\[
\pi(r_1+r_2)\Delta s
\]
である。円すい台の場合は，大きい円すいと切り取った円すいの側面積の差を，相似比で整理するとこの式を得る。円筒面と円環の場合も，それぞれの面積公式から確かめられる。線分を十分短くすれば，$r_1,r_2$ はほぼ同じ $r$ なので，これは $2\pi r\Delta s$ になる。

回転軸と同じ平面内の曲線が区分的に連続微分可能で，回転によって同じ面を重複して作らないとき，折れ線近似の極限として回転面の面積は
\begin{equation}\label{式:回転面の面積}
S=2\pi\int_C r\dd s
\end{equation}
になる。ここでは，この極限の正当化の詳細は省く。たとえば $y=f(x)\geq0$ を $x$ 軸のまわりに回すなら
\[
S=2\pi\int_a^b f(x)\sqrt{1+f'(x)^2}\dd x.
\]
ここで得られるのは，指定した曲線が作る面の面積である。立体の両端を塞ぐ円板などは，その曲線に含めていなければ別に数える。

\begin{theorem}{パップス＝ギュルダンの面積に関する定理}
一様な平面曲線の長さを $L$，その線重心と，曲線と同じ平面内にある回転軸との距離を $R$ とする。曲線が軸の片側にあり，一回転によって同じ面を重複して作らない場合，曲線が作る回転面の面積は
\[
S=2\pi R L
\]
である。
\end{theorem}

実際，軸を $x$ 軸，曲線を $y\geq0$ に取れば $r=y$ であり，線重心の定義から $\int_C y\dd s=L\overline y=LR$ となるので，式\eqref{式:回転面の面積}から従う。回転軸の位置や方向が違っていても，軸に合わせて座標を取り直せば同じ議論になる。

二つの定理で掛ける円周は，図形の端や幾何学的な中心の円周ではなく，それぞれの重心が進む円周である。また，体積の定理には面積 $A$ と面積重心，面の定理には長さ $L$ と線重心を対応させる。

\begin{example}{半円板を外側の軸のまわりに回す}
$a>0$，$b>0$ とする。上半円板 $x^2+y^2\leq a^2$，$y\geq0$ を，直線 $y=-b$ のまわりに一回転してできる立体の体積と全表面積を求めよ。
\end{example}

図\ref{図:半円の二つの重心}のように，回転軸は図形の外にあり，重心の高さに $b$ を加えたものが回転半径になる。半円板の面積重心を使うと，
\[
\begin{aligned}
V&=2\pi\left(b+\frac{4a}{3\pi}\right)\frac{\pi a^2}{2}\\
 &=\pi^2a^2b+\frac{4\pi a^3}{3}.
\end{aligned}
\]
断面で確かめるなら，外半径 $b+\sqrt{a^2-x^2}$，内半径 $b$ の円環である。したがって
\[
\begin{aligned}
V&=\pi\int_{-a}^a
\left\{\left(b+\sqrt{a^2-x^2}\right)^2-b^2\right\}\dd x\\
 &=\pi\int_{-a}^a(a^2-x^2)\dd x
  +2\pi b\int_{-a}^a\sqrt{a^2-x^2}\dd x,
\end{aligned}
\]
となり，第一項が $\frac{4\pi a^3}{3}$，第二項が $\pi^2a^2b$ で一致する。

\begin{figure}[htbp]
\centering
\begin{tikzpicture}[scale=1.25,>=Stealth]
\fill[blue!10](-2,0)arc[start angle=180,end angle=0,radius=2]--cycle;
\draw[blue!75!black,thick](2,0)arc[start angle=0,end angle=180,radius=2];
\draw[thick](-2,0)--(2,0);
\draw[->](-2.5,0)--(2.6,0)node[right]{$x$};
\draw[->](0,-1.5)--(0,2.45)node[above]{$y$};
\draw[orange!80!black,thick](-2.5,-1)--(2.5,-1)node[right]{$y=-b$};
\fill[blue!75!black](0,0.8488)circle(0.045);
\fill[red!75!black](0,1.2732)circle(0.045);
\draw[blue!75!black](0,0.8488)--(2.5,0.6)node[right]{面積重心 $\frac{4a}{3\pi}$};
\draw[red!75!black](0,1.2732)--(2.5,1.6)node[right]{線重心 $\frac{2a}{\pi}$};
\draw[<->](-2.25,-1)--(-2.25,0)node[midway,left]{$b$};
\node[below]at(-2,0){$-a$};\node[below]at(2,0){$a$};
\end{tikzpicture}
\caption{同じ半円でも，板と円弧では重心が違う。回転半径は軸から測る。}\label{図:半円の二つの重心}
\end{figure}

表面積では，境界を円弧と直径に分ける。円弧が作る面は，円弧の長さ $\pi a$ と線重心の高さ $\frac{2a}{\pi}$ を使って
\[
S_{\mathrm{円弧}}
 =2\pi\left(b+\frac{2a}{\pi}\right)\pi a
 =2\pi^2ab+4\pi a^2.
\]
これだけでは全表面積にならない。直径 $-a\leq x\leq a$，$y=0$ も，回転すると半径 $b$，軸方向の長さ $2a$ の円筒面を作る。その面積は
\[
S_{\mathrm{直径}}=2\pi b\cdot2a=4\pi ab.
\]
両者は端の円周で接するが，面積のある部分では重ならない。よって全表面積は
\[
S=2\pi^2ab+4\pi a^2+4\pi ab.
\]
$x=\pm a$ では外半径と内半径がともに $b$ となり，厚みが零なので，さらに端の円板を加える必要はない。どの境界がどの面を作るかを確かめた結果として，全表面積が決まる。

面積，体積，長さ，道のり，平均はいずれも積分で表せるが，同じ被積分関数を使うわけではない。細い帯なら高さと幅，薄い立体なら断面積と厚さ，短い曲線なら二方向の変化から作った距離が必要になる。公式を選ぶ前に，細かく分けた一つ分が何を表すかを言えるようにすると，符号や重複の誤りに気づきやすくなる。

